% Experiments
\section{Experiments and diagnostics}
\label{sec:experiments}

We run twelve experiments to test \ISR against the global census baseline.
Table~\ref{tab:kernel_roles} summarizes kernel roles.

\begin{table}[t]
  \centering
  \caption{Kernel roles, evidence status, and rationale. Evidence kernels are used in main ISR results; diagnostic kernels illustrate sensitivity to design choices.}
  \label{tab:kernel_roles}
  \begin{tabular}{lllp{5cm}}
    \toprule
    Kernel & Role & Evidence & Reason \\
    \midrule
    Gaussian & Primary locality kernel & Evidence & Exponential decay in field-space distance; standard smoothing kernel. \\
    Exponential & Robustness kernel & Evidence & Heavier tail than Gaussian; tests sensitivity to kernel shape. \\
    \texttt{same\_basin} & Diagnostic & Diagnostic & Trivially conditions on $x_0$ basin membership; no field-distance resolution. \\
    \texttt{ancestry\_proxy} & Diagnostic & Diagnostic & Uses trajectory index as a proxy for branch ancestry; not a true causal ancestry graph. \\
    \bottomrule
  \end{tabular}
\end{table}


\subsection{Experiment 1: Global census baseline}
\label{exp:1}

Compute $P_{\rm global}(v)$ at each cutoff $N_k$.
\emph{Purpose}: quantify cutoff dependence of naive counting.
\emph{Expected failure}: distribution shifts toward the largest basin as $N$ increases.

The global census is the natural null hypothesis: it treats all generated sectors
as equally relevant regardless of their field-space location.
A measure must justify why some sectors count more than others, and \ISR provides
a specific locality-based answer.
The global census is therefore the correct baseline---not because it is expected
to work well (it is not), but because it represents the simplest possible measure
that makes no distinction between sectors.
If \ISR does not improve upon it, the locality proposal adds no value.

\subsection{Experiment 2: ISR locality-weighted measure}
\label{exp:2}

Compute $P_{\rm ISR}(v \mid x_0)$ using a Gaussian kernel ($\ell = 1$)
with field-space distance.
\emph{Purpose}: test whether locality weighting shifts probability toward the
$x_0$ basin and improves stability.

\subsection{Experiment 3: Cutoff stability}
\label{exp:3}

Compare $D_{\rm KL}$ and Wasserstein distances across all cutoffs
for both measures, and for multiple kernels.
\emph{Purpose}: identify which kernel---measure combinations converge.

\subsection{Experiment 4: Kernel sensitivity}
\label{exp:4}

Vary $\ell$ and kernel functional form to check sensitivity.
\emph{Purpose}: ensure results are not an artifact of one kernel choice.

\subsection{Experiment 5: Basin-boundary diagnostics}
\label{exp:5}

At fixed distance thresholds, compute the fraction of ``nearby'' sectors
($d < d_{\rm thr}$) that lie in a different basin from $x_0$.
\emph{Purpose}: measure how often the locality assumption fails.

\subsection{Experiment 6: Configuration-space locality diagnostics}
\label{exp:6}

For three distance metrics (field-space, initial-field, basin-transition),
evaluate continuity score, basin-exception rate, and locality-failure rate.
\emph{Purpose}: characterize the structure of locality violations.

\subsection{Experiment 7: Kernel-weighted sector averaging}
\label{exp:7}

Compute $c_{\rm eff}(\ell)$ at varying $\ell$ to show unobserved-sector mixing.
\emph{Purpose}: illustrate kernel-weighted sector averaging.

\subsection{Experiment 8: Stress test}
\label{exp:8}

Repeat experiments~1 and~2 for $\sigma \in \{0.5, 1.0, 1.5, 2.0, 3.0\}$ with
seeds $1$ through $20$ ($100$ total runs), $N = 2,\!000$ per run.
\emph{Purpose}: test robustness across random and systematic variation.

\subsection{Experiment 9: Ell sensitivity scan}
\label{exp:9}

Vary $\ell$ across 15 values in logspace from 0.05 to 10.0 with Gaussian and
exponential kernels.
For each $\ell$, record final KL drift, effective sample size $N_{\rm eff}$,
basin exception rate, effective coefficient $\Lambda_{\rm eff}$, and
convergence rate.
The convergence rate is defined as the slope of log KL drift versus log cutoff,
$d(\log{\rm KL}) / d(\log N)$, estimated by linear regression over the
log--log curve.
A convergence rate of $-1$ indicates ${\rm KL} \propto 1/N$, matching the
parametric rate of a well-behaved Monte Carlo estimator; rates approaching $0$
indicate that additional samples do not improve the estimate.
Also run a randomized-distance null control (shuffled sector--distance
pairings) to verify that the improvement is not an artifact of the kernel
smoothing.
\emph{Purpose}: map the sweet spot of $\ell$ where ISR provides meaningful
improvement without being trivially narrow or approaching global census.

\subsection{Experiment 10: Sigma--ell boundary}
\label{exp:10}

Grid over $\sigma \in \{0.5, 1.0, 1.5, 2.0, 3.0\}$ and
$\ell \in \{0.1, 0.25, 0.5, 1.0, 2.0\}$, 5 seeds per cell
($125$ total runs), $N = 2,\!000$ per run.
\emph{Purpose}: identify any $(\sigma, \ell)$ region where ISR underperforms
the global census.

\subsection{Experiment 11: Flat-parent composition sanity check}
\label{exp:11}

For each cutoff, compose the \ISR weighting with a flat parent measure
(equal weight to all sectors). The composition applies \ISR re-weighting
on top of the parent baseline and checks that the resulting
cutoff-stability trend is consistent with the standalone \ISR result.
\emph{Purpose}: sanity check that \ISR can be stacked with a trivial
parent measure without introducing numerical inconsistencies.

\subsection{Experiment 12: Landscape perturbation robustness}
\label{exp:12}

Shift all well centers $\mu \to \mu + 0.5$ in the toy landscape and repeat
Experiments~1 and~2 at $\sigma \in \{1.5, 3.0\}$, 5 seeds per cell.
The perturbed landscape preserves the same number of basins and the same 
basin physics vectors $\boldsymbol{\theta}$, but shifts basin boundaries
relative to $\phi_0$.
\emph{Purpose}: verify that \ISR improvement is not an artifact of the
specific well placement in the unperturbed landscape.
